\subsection{The matched local data from the given period seed}
\label{sec:flag-local-data}

The preceding diagonal estimate retains the original
\(\mathrm{LocDatC}\) data. We now supply them in the concrete
standing regime in which the amplitude theorem uses them.
In particular we use the given
\(\mathrm{ConcTie}\), including integrability of \(X.\mathrm{ph}\),
and the given arithmetic ties, including \(\mathrm{dsigTie}\).
We do not require the stronger universal \(\mathrm{BesselOfW}\)
statement without the integrable-projector hypothesis.

\begin{lemma}[Extraction by a discrete-series projector]
\label{lem:flag-bessel-extraction}
Let \(G\) be a real reductive group, \((\sigma,\mathcal H)\)
an irreducible unitary square-integrable representation with
formal degree \(d>0\) for Haar measure \(dg\), and \(w\) a
smooth unit vector. Let \(S\) be a closed unimodular subgroup
and \(\xi:S\to\mathrm U(V)\) a smooth unitary representation
on a finite-dimensional Hilbert space. Let \(W:G\to V\) be
continuous, satisfy \(W(sx)=\xi(s)W(x)\), and have nonzero
finite squared norm for the invariant quotient measure.
Suppose
\[
 e(g)=d\,\overline{\langle w,\sigma(g)w\rangle}\in L^1(G),
 \qquad R(e)W=W ,
 \tag{LD1}
\]
where \(R(g)W(x)=W(xg)\) and \(R(e)\) is the integrated
operator on the equivariant quotient Hilbert space.
Then there is a linear map \(B:\mathcal H^\infty\to V\)
satisfying
\[
 B(\sigma(s)v)=\xi(s)B(v),\qquad
 W(g)=B(\sigma(g)w),
\]
and bounded on smooth vectors by a finite sum of the
Hilbert norms of enveloping-algebra derivatives.
\end{lemma}
\begin{proof}
The Hilbert space of equivariant \(L^2\)-sections has a strongly
continuous unitary right action. Schur orthogonality for the
square-integrable representation gives, in the group
convolution algebra with the stated Haar measure,
\[
 e^*=e,\qquad e*e=e,\qquad
 e*\delta_g*e=\langle w,\sigma(g)w\rangle e.
 \tag{LD2}
\]
For the last identity, at \(x\) the integral on the left is
\[
 d^2\int_G
 \overline{\langle w,\sigma(y)w\rangle}
 \langle \sigma(x)w,\sigma(y)\sigma(g)w\rangle\,dy
   =\langle w,\sigma(g)w\rangle e(x),
\]
using the conjugate-linear first-argument convention.
The convolution is defined in \(L^1\); the displayed integral
also converges by Cauchy--Schwarz and square integrability.
The same calculation at \(g=1\) gives idempotence.
Thus \(R(e)\) is an orthogonal projection and (LD1)--(LD2)
give
\[
 \langle R(g_1)W,R(g_2)W\rangle
   =\|W\|^2\langle \sigma(g_1)w,\sigma(g_2)w\rangle .
\]
It follows that
\(\sigma(g)w\mapsto R(g)W\) extends linearly to a
well-defined scaled isometry \(T:\mathcal H\to L^2(S\backslash
G;\xi)\), because the orbit span of \(w\) is dense by
irreducibility. The isometry intertwines the two actions.

For clarity, evaluation is made on smooth vectors, not on
arbitrary \(L^2\) representatives. On the smooth homogeneous
space \(S\backslash G\) with its associated smooth vector bundle,
a smooth vector of the right action is a smooth section.
Indeed, on a relatively compact local trivialization the
right-action derivatives are differential operators with
smooth coefficients, whose tangent parts span the tangent
space. Locally any coordinate derivative is a smooth
linear combination of these and a zeroth-order operator.
Applying this repeatedly and multiplying by a compactly
supported cutoff bounds ordinary Sobolev norms of any fixed
order by finitely many quotient \(L^2\) norms of right-action
derivatives. The invariant quotient measure has smooth positive
density on the chart. Sobolev embedding, first in arbitrarily
high orders and then in order strictly above half the chart
dimension, proves both smoothness and bounded evaluation.
This argument applies to weak derivatives of \(L^2\)-sections:
smoothness of the orbit map gives all such derivatives in
\(L^2\). In particular for a real Lie-algebra basis there
exist a finite set \(\mathcal D\) of words in that basis
(including the empty word) and a finite constant \(C\) such
that
\[
 \|(Tv)(1)\|\le C\sum_{D\in\mathcal D}
            \|R_*(D)Tv\|_{L^2}
       = C\|W\|\sum_{D\in\mathcal D}\|\sigma_*(D)v\|.
\]
Set \(B(v)=(Tv)(1)\) using the unique smooth section.
The intertwining property and equivariance of the section
give the first required identity. Since \(Tw=W\) in \(L^2\),
the smooth representative agrees everywhere with the
given continuous \(W\), by the full support of the quotient
measure. This proves the second identity and the asserted
derivative bound. The bound is qualitative, as required
in \(\mathrm{IsCWCont}\); no rank-uniform Sobolev constant
is used.
\end{proof}

\begin{theorem}[The actual local-data consumer]
\label{thm:flag-local-data}
For the concrete setting used in the original amplitude
theorem, with its given \(\mathrm{ConcTie}\) and
\(\mathrm{dsigTie}\), there exist the original
\(C:\mathrm{Choice}\ X\ 1\) and \(h_B:\mathrm{WuBessel}\ C\)
with \(\mathrm{GoodLoc}(C,h_B)\).
The assertion is simultaneous for the actual words and real
places. In particular the local-data hypothesis of (FL8)
does not require a further local existence conjecture.
\end{theorem}
\begin{proof}
We give the construction using the original frames and measures.
At a real place the two given isotropic frames form dual
hyperbolic bases. Complete their scaled vectors to a Hermitian
basis of the given signature by taking an orthonormal basis
of the orthogonal complement. The first frame is scaled
by \(\ell^{r-a}\) and the dual by
\(\overline{\ell}^{-(r-a)}\), where
\(\ell=\mathrm{frS}(D)=\mathrm{ka}/(2\pi)\ne0\).
This yields exactly the matrices \(Q_\tau,Q_\tau^{-1}\)
of \(\mathrm{IsTranspC}\), independently of the word.
It transports each given \(z_u\) into the middle space
and preserves its norm \(-\mathrm{Tr}(D,u)\), which is
strictly negative. The compact stabilizer of \(z_u\)
in this middle space is the unitary group of its positive
orthogonal complement. Transport of the given orthonormal
frame \(wfr_{u,\tau}\) identifies its prescribed compact
representation with \(\mathrm{piB}\); its matrix entries
are precisely those of the original \(h_J\)-pairings.
Thus the column highest weight is unchanged.

The transported discrete series is \(\mathrm{sigB}\).
The two Hermitian signature models are identified by the
given signature ordering and a Cayley isometry to \(Hpk\).
Conjugating through that isometry identifies the
Harish--Chandra parameter with \(\mathrm{lamD}(D,\tau)\):
the original \(\mathrm{eps\_eq}\) and
\(\mathrm{lam\_eq}\) fields give these equalities entrywise.
In particular it has the required discrete-series parameter.
The original \(w_\tau\) is a unit vector in the lowest
compact type. The maximal compact subgroup transported
by \(Q_\tau\) is conjugate to the compact subgroup in the
chosen Cayley model. Choose \(y_\tau\) realizing that
conjugation with the same signature ordering.
Equivalently, if \(A=Q_\tau P^{-1}\) in signature
coordinates (using the given sign-coordinate identification),
the Cayley-conjugated representation is
\(g\mapsto\sigma(A g A^{-1})\), and one may take
\(y_\tau=Q_\tau^{-1}A Q_\tau\).
Then \(v_{0,\tau}=\mathrm{sigB}(y_\tau)w_\tau\) is
exactly a unit lowest-type vector for this representation.
Both \(w_\tau\) and \(v_{0,\tau}\) are smooth. We use the
classical Harish--Chandra admissibility theorem for
irreducible unitary representations of real reductive
groups and its consequence that their compact-finite
vectors are smooth (Knapp, \emph{Representation Theory
of Semisimple Groups}, Chapter VIII). Concretely, in
each finite-dimensional compact isotypic space the
smooth vectors are dense by convolution smoothing and
compact projection, hence equal the whole space.
The lowest-type subspace consists of compact-finite
vectors. This is the precise smoothness needed here.

Use the pushforward \(\mu_{G,\tau}\) of the original real
Haar measure by \(Q_\tau^{-1}(\cdot)Q_\tau\), probability Haar
measure \(\mu_{W,u,\tau}\) on the compact stabilizer of the
last vector, and exactly the original
\(\mu_{N,\tau}=\mathrm{haarLie}(N,\mathrm{kapN}(D,1))\).
The latter is positive Haar measure by the unipotent
coordinate calculation, with the specified positive scalar.
The full stabilizer \(S_{u,\tau}\) is the semidirect
product of this compact group and the flag unipotent.
Indeed reduction on the isotropic quotient fixes \(z_u\)
and acts on its orthogonal complement, and the kernel is
the unipotent radical; the fixed dual hyperbolic frame
gives the splitting. Conjugation by the compact factor
preserves unipotent Haar measure (its modular multiplier
is a positive character of a compact group).
Thus the stated product \(\mu_S\) is Haar, and \(S_{u,\tau}\)
is unimodular and closed. Choose an invariant quotient
measure by the Weil formula with these measures.
These are the data of \(\mathrm{Choice}\), including all
\(\mathrm{GoodW}\) measure and norm conditions.
The choices of \(Q_\tau,y_\tau\) were common to the words;
the compact stabilizer and quotient measure use the actual
word vector.

Write \(W_{u,\tau}\) for the given \(\mathrm{Wloc}\) in
these coordinates, and \(e_\tau\) for its prescribed
formal-degree-normalized matrix coefficient, as in
\(\mathrm{eFun}\). We verify the local hypotheses directly
from the actual global seed. The given pure-tensor map
is the restriction of a unitary tensor identification:
the identity on inner products makes the linear extension
well-defined and isometric, and its span condition makes
it onto. The global squared norm is one, so no local
factor vanishes identically. Apply global equivariance to
an element supported at just one real place and to fixed
nonzero values at all the other places. Cancellation of
their nonzero tensors gives the local compact and
unipotent equivariance. The character is the original
one after coordinate transport: each root coordinate
\(\langle f_a,\kappa\,\mathrm{nxt}(e,z)_a\rangle\)
becomes \(\ell\) times the corresponding normalized root
coordinate, including the final root evaluated on \(z_u\).
Hence their sum becomes \(\ell L\), where \(L\) is the
sum in the local character. The actual exponential has
argument \(2\pi i(\ell L+\overline{\ell L})/g\),
which is \(i(2\,\mathrm{ka}/g)\operatorname{Re}L\).
This is precisely \(\mathrm{psiS}\), not a freely
rescaled character. The compact action uses the given
tensor equivariance and \(wfr\)-basis. These verify the
local \(\mathrm{IsSPsiEquiv}\) assertion.

The global unipotent and compact measures differ from
the products of the corresponding chosen local measures
only by strictly positive finite Haar constants.
The Weil formulas then identify the global quotient
measure on invariant functions with the inverse scalar
times the product local quotient measures. Consequently
the actual global norm-one condition, by Tonelli, implies
that every \(\|W_{u,\tau}\|_{L^2(S_{u,\tau}\backslash G_\tau)}^2\)
is positive and finite: all factors are positive since
the local functions are continuous and nonzero, and a
single infinite factor would make the product infinite.
All measures have full support on the quotient.
This argument keeps the original \(W_{u,\tau}\), without
making separate choices of their scalar normalizations.

By the given formal degree definition and Schur
orthogonality, \(\mathrm{ph}\) is precisely the product of
the \(e_\tau\)'s, for the product real Haar measure
(\(\mathrm{phInf\_eq}\)). Each coefficient is nonzero
near the identity. Integrability of the actual global
\(\mathrm{ph}\), supplied by \(\mathrm{ConcTie}\), and
Tonelli therefore give \(e_\tau\in L^1\) at every place.
The original global pointwise convolution identity uses
this product kernel. In the quotient Hilbert space it
implies
\[
     \bigotimes_\tau R(e_\tau)W_{u,\tau}
                =\bigotimes_\tau W_{u,\tau}.
\]
Indeed the integrated operator is bounded by the
\(L^1\)-norm of the kernel, the product right action is
the action on the tensor product of quotient-section
spaces, and Fubini gives the displayed factorization.
The global pointwise identity agrees almost everywhere
with the integrated operator: on relatively compact
quotient charts, integration against compactly supported
test sections and Fubini identify the two, since the
integrated \(L^2\) norm is finite. No convergence beyond
the given \(L^1\) kernel and \(L^2\) sections is required
for that identification. By (LD2) each \(R(e_\tau)\)
is an orthogonal projection. Equality of the pure-tensor
norms and contraction in every factor of positive norm
give \(R(e_\tau)W_{u,\tau}=W_{u,\tau}\) in its own
Hilbert space.

The original local-data definition also requires the
pointwise convolution identity and integrability at every
point. These follow from the given global pointwise
assertion, not just from equality in \(L^2\).
Fix one place and any point there; choose points at the
other places where their local \(W\)'s are nonzero.
Continuity and the nonzero coefficient at identity
give positive integrals of their absolute convolution
integrands on small neighborhoods. The global absolute
integrability and Tonelli give absolute integrability
of the local convolution at the chosen point.
The local convolution functions thus exist everywhere.
Fubini in the global reproduction identity gives equality
of the two products of tensors at all tuples of points.
Fix nonzero anchor values in the factors; this equality
forces the local convolution functions to be nonzero
constant multiples of their respective original functions
(with constants independent of the point in a factor).
Their \(L^2\) identities just proved make each multiple
one. Thus the pointwise identities hold everywhere,
as required in \(\mathrm{LocFactor}\).

Apply Lemma~\ref{lem:flag-bessel-extraction} in each
transported local group to the given \(W_{u,\tau}\).
It constructs a map \(B_{u,\tau}\) on the smooth vectors,
whose restriction to the given orbit is precisely the
given \(\mathrm{Wloc}\). Extend it by zero off the smooth
vectors as a function on the Hilbert space. The original
\(\mathrm{IsBesselMapSm}\) requires linearity and
equivariance only on smooth vectors, so this extension
satisfies its exact domain requirements. The derivative
bound in that lemma is the original \(\mathrm{IsCWCont}\)
with a finite list of Lie-algebra words. Smoothness on
the real matrix group agrees with the specified
\(\mathrm{SmoothOnH}\), since the group is an embedded
smooth submanifold and a smooth map on it extends
locally by a smooth chart retraction.
We obtain \(\mathrm{WuBessel}\ C\) simultaneously,
without a Bessel multiplicity-one assertion.

Finally \(\Phi_{u,\tau}(x)=B_{u,\tau}(\mathrm{sigB}(x)
v_{0,\tau})=W_{u,\tau}(x y_\tau)\) in the normalized
coordinates. Right translation is unitary on the
quotient Hilbert space. Its squared norm is therefore
the very same positive finite local mass. Define
\(c_N\) as this squared norm divided by the actual compact
dimension, as in \(\mathrm{cOf}\); it is strictly positive.
The diagonal scalar \(\mathrm{sclOf}\) equals one, and
the two squared-norm requirements of \(\mathrm{GoodS}\)
are now identities. For the formal degree requirement,
in the given Cayley model the transported discrete series
has exactly the given signature and Harish--Chandra
parameter. Schur orthogonality for its Haar measure
scaled by the original \(\mathrm{cdScale}\) provides
its formal degree; applying the given \(\mathrm{dsigTie}\)
identifies that degree with \(D.\mathrm{dsig}(\tau)\).
This is the required formal degree for the scaled
measure, distinct from the unscaled \(d\) in (LD1).
All remaining conditions of \(\mathrm{GoodS}\) have been
verified above. This proves \(\mathrm{GoodLoc}(C,h_B)\).
\end{proof}

\begin{corollary}[Diagonal zero-Gram bound in the concrete standing regime]
\label{cor:flag-zero-gram-concrete}
Under the original amplitude hypotheses (including the given
\(\mathrm{Ties}\), \(\mathrm{KTwoHypotheses}\) and
\(\mathrm{ConcTie}\)), and the original absolutely convergent
orbital expansion, the diagonal nonidentity zero-Gram
operator satisfies
\[
 \sum_{t\in\mathrm{degSet}}\|\mathrm{op}(t)\|_{\rm HS}<\infty,
 \qquad
 \|K''_{uu}\|_{\rm HS}^2\le d_u e^{-kf\log k}.
\]
This is the actual diagonal zero-Gram norm estimate needed
on the concrete data. It leaves the generic-return estimates
and the subsequent arithmetic comparison open.
\end{corollary}
\begin{proof}
Theorem~\ref{thm:flag-local-data} supplies the matched data of
Theorem~\ref{thm:flag-zero-gram-saving}. The estimates and
thresholds in (GA1)--(GA8) are unchanged; the construction
here adds no numerical threshold. Absolute convergence of
the restricted operator series follows by restricting the
given \(\mathrm{AbsConv}\) series. Apply (FL8).
\end{proof}
